% !TeX root = ../papers/nc-p3.tex
\subsection{接空間が三次元の有限点からのOre表示}
\label{sec:maximal-tangent-ore}
有限四辺で最も大きい接空間を持つ点を用いると、三生成部分の存在を仮定せずに
四生成の正則性を三生成の正則性へ還元できる。
以下の周期に関する点の同一視にはFrobenius標準四窓を使う。

\begin{theorem}\label{thm:maximal-tangent-ore}
標準有限Frobenius候補の四辺が全て有限で、ある辺の閉点の接空間が三次元とする。
その点から両方向に延びる点列を $[p_i]\in\mathbb P(V_i^*)$ とする。
\begin{enumerate}
\item $[p_{i+4}]=[p_i]$ であり、このような列は始点を固定すると高々二本である。
$H_i=\ker p_i$ によって内在的に三生成部分が定まる。
\item 各位置で一次基底を選ぶと、ある一定の三生成三関係の二次代数
$B=T(H)/(S)$ と写像 $d_i:H\to B_2$ により、候補の全関係は
\[
         S=0,\qquad X_i h-h X_{i+1}=d_i(h)\quad(h\in H)
\]
になる。ここで隣接する $H$ は選んだ同一視で表している。三次における条件は
\begin{equation}\label{eq:ternary-shift-derivation}
       (d_i\otimes 1+1\otimes d_{i+1})(S)=0\quad\text{in } B_3
\end{equation}
で、この候補では全始点で成立する。各区間の長さ $n$ の空間は
$\bigoplus_{r=0}^n B_{n-r}X^r$ と同型である。特に $B$ は実際に埋め込まれる。
\item 候補が標準14641型Koszul左右AS正則であることと、$B$ がHilbert級数
$(1-t)^{-3}$ の三次元quadratic Koszul左右AS正則代数であることは同値である。
逆構成ではこのような $B$ と任意の列 $d_i$ について
\eqref{eq:ternary-shift-derivation}だけで標準AS正則性を得る。
この逆構成全体には辺が正次元のものも含まれる。
\end{enumerate}
\end{theorem}
\begin{proof}
系\ref{cor:finite-edge-point-sequences}から点列は一意で、全位置で同じ三次元接空間を持つ。
$X_i$ を $p_i(X_i)=1$ と選ぶ。関係の $X_iX_{i+1}$ 成分は零である。
$R_i$ から $X_iH_{i+1}$ と $H_i X_{i+1}$ の二成分への写像はそれぞれ階数3で、
両方を合わせた写像は六つの双一次式の点 $(p_i,p_{i+1})$ における
Jacobi行列の転置である。接空間次元が3なので合わせても階数3である。
従って両核は同じ三次元空間 $S_i=R_i\cap(H_i\otimes H_{i+1})$ で、残る関係は
\[
 X_i h=f_i(h)X_{i+1}+d_i(h),\qquad f_i:H_{i+1}\xrightarrow{\sim}H_i
\]
と書ける。補正は $H_i\otimes H_{i+1}$ の $S_i$ を法とした類を使う。

まず $B^{(0)}$ を $H_i,S_i$ だけで定義したindexed代数とし、
$d=\dim((S_i\otimes H_{i+2})\cap(H_i\otimes S_{i+1}))$ とする。
純粋な三次の商次元は $27-18+d=9+d$ である。
混合語の $X$ を右へ動かせば、四生成の三次全体は
$(9+d)+6+3+1=19+d$ 個で生成される。
四生成の三次交叉 $W_{i,3}$ の次元4から三次商は20次元なので $d\ge1$ である。

一方、純粋な交叉は $W_{i,3}$ のうち第一位置を $p_i$ で縮約して零になる部分に入る。
定理\ref{thm:frobenius-tensor-reconstruction}により $W_{i,3}$ は四テンソルの
第四位置の縮約と同一視される。この写像に第一位置の $p_i$ を続けたものの核は
$R_{i-1}(p_i)$ の零化空間であり、直前の点 $kp_{i-1}$ である。
ここで第四位置は標準四周期性によって位置 $i-1$ と読んだ。
よって $d=1$ で、純粋な交叉の生成元は $w(p_{i-1})$ である。
その末端支持は次元3で、直前の辺の零点対から $\ker p_{i-2}$ に等しい。
純粋性から同じ支持は $H_{i+2}$ に含まれ、従って $[p_{i+2}]=[p_{i-2}]$ となる。
先頭支持も次元3である（命題\ref{prop:three-finite-edges}の証明）。
また各局所点環は埋込み次元3の完全交叉なので、次数2の次数商の次元は
$6-3=3$ 以上である。三乗イデアルが零ならこの商がsocleに入り
Gorenstein性に反する。従って長さは $1+3+3+1=8$ 以上で、全長20から個数を得る。

上述の20個の三次生成元は基底である。$X_i S_{i+1}$ を右へ還元すれば、
一つの $X$ の係数は $(f_i\otimes f_{i+1})S_{i+1}$ である。
還元結果が零なので $S_i$ に含まれ、両方三次元で等しい。
全整数上で $f_i$ に沿って順次 $H_i$ を同一視すれば、$S_i=S$、
$f_i=1$ となる。この基底に四周期性を要求していない。
$X$ を含まない残りの係数がちょうど\eqref{eq:ternary-shift-derivation}である。

この三次条件で高次に新しい障害がないことを直接示す。
$d_i$ を積の規則 $D_i(ab)=D_i(a)b+a D_{i+\deg a}(b)$、
$D_i(1)=0$ により自由代数から $B$ へ延ばす。
全始点で関係 $S$ を殺すので、その両側イデアルも殺し $B$ 上に降りる。
位置ごとに $bZ^r$ の形式的直和を作り、左から $h$ は係数に掛け、
$X_i$ は $bZ^{r+1}+D_i(b)Z^r$ に送る。積の規則によって全関係が作用で零になる。
右へ還元した語を終点の $1$ に作用させるとその係数を復元できるので、
表示した直和が全次数で成立する。$X$ の個数で濾過した随伴次数は
中心変数 $Z$ 一個を加えた $\check{B[Z]}$ となる。

$B$ が定理の三次元正則代数なら、$B[Z]$ は二つの線型解消のテンソル積によって
四次元Koszulで、二次双対は $B^!$ と一生成外積代数の次数付きテンソル積で
14641型の完全対合を持つ。元の候補で各内部次数のbar複体を $X$ の総数で濾過する。
濾過は有限で、その随伴次数は直前の代数のbar複体である。
そのホモロジーは内部次数とホモロジー次数が等しいところだけにあるので
スペクトル列はそこで退化し、元の代数もKoszulである。
barの分割余積が濾過と両立するから、有限次元の対角部分を双対化して得る
Extの随伴次数の積も上記の二次双対の積になる。
従って元の二次双対にも同じ次元と完全対合が持ち上がり、左右AS条件を得る。
Hilbert級数は直和表示から $(1-t)^{-4}$ である。

逆に元の代数 $A$ が標準ASなら直和表示から $H_B(t)=(1-t)^{-3}$ である。
$B$ の右自明加群の極小次数付き自由解消を、始点を固定して
$\check B\to A$ に沿って誘導する。各終点では左 $B$ 側で上の直和を取っているので
誘導は完全であり、微分の正次数性も保たれる。
得られるのは右点加群（各長さで $kX^r$）の極小解消である。
$A$ の大域次元4により $g=\operatorname{pd}_B k\le4$ となる。
以下、無限個の生成元を持つ次数付き加群にも使える形で、$B$ の両側のsocleが零であることを確認する。
左右の $k$ の極小解消の最大長は $\operatorname{Tor}^B(k,k)$ の最終次数で一致する。
従って任意の下に有界な次数付き右加群 $M$ に対して
$\operatorname{Tor}^B_{>g}(M,k)=0$ であり、左加群でも同様である。
正の大域次元 $g$ で同次socle元があれば $k(-a)\hookrightarrow B$ があり、
その余核の $\operatorname{Tor}_{g+1}$ は $\operatorname{Tor}_g(k(-a),k)\ne0$ に同型となる。
これは矛盾である。

$B$ の生成元、三関係、唯一の純粋な三次交叉から右Koszul複体の先頭
$B(-3)\to B(-2)^3\to B(-1)^3\to B\to k$ を作る。
最左写像の三係数は交叉生成元の末端支持が全三次元であるため一次生成元の基底である。
従って左socleが零であることから単射である。
生成関係の定義から $B(-1)^3$ 以降は完全で、残る一箇所のホモロジーは
既知のHilbert級数とEuler等式から零である。逆側も先頭支持を使って同じ結論を得る。
両端の三次元支持は $B^!$ の次数1と2の両向きの積を完全にする。
以上で $B$ は三次元quadratic Koszul左右AS正則である。
\end{proof}

\begin{corollary}\label{cor:maximal-tangent-ore-isomorphism}
上記の点列を指定した二対象の正規表示を $(B,d_i),(B',d'_i)$ とする。
指定列を保つ添字保持同型の存在は、次数付き代数同型 $G:B\to B'$ と
各位置の $v_i\in B'_1$ が存在して
\[
       Gd_i(h)=d'_i(Gh)+v_i(Gh)-(Gh)v_{i+1}\quad(h\in H)
\]
となることと同値である。列を指定しない場合は、始点の指定点の像として
標的の高々二つの三次元接点を試せば十分である。
\end{corollary}
\begin{proof}
指定点を保つ一次同型は $h\mapsto g_i h$, $X_i\mapsto c_i X'_i+u_i$
の形である。$c_i\ne0$ で、混合関係の $X'$ 係数から
$c_i g_{i+1}=c_{i+1}g_i$ となる。
各位置の一次写像を $c_i$ で割っても二次関係の像は同じである。
従って $g_i/c_i$ は一定の $G$ となり、純粋な三関係を対応させる。
残る係数が表示した式である。逆にこれらの式は全関係を対応させ、一次写像が可逆なので同型を与える。
\end{proof}

\path{computations/verify_maximal_tangent_ore.py} では、
四式 $x^2+4yz+2wx,y^2+4xz,z^2+4xy,w^2$ の対称AS窓で辺の有限性を検査する。
平方点での接空間、純粋な三関係、三次交叉の両端支持、四生成の三次元数を有理数で計算する。
一つの位置の混合補正だけを変えると、同じ接空間と純粋関係を保っても三次元数が20未満になる。

\begin{theorem}[一位置の線型パラメータと移動作用]\label{thm:ternary-ore-linear-space}
三次元Koszul左右AS正則な $B=T(H)/(S)$、$H_B(t)=(1-t)^{-3}$ を固定する。
$D=\operatorname{Hom}(H,B_2)$ は18次元で、$M=\operatorname{Hom}(S,B_3)$ に値を取る写像
$\lambda d=(d\otimes1)|_S$, $\rho d=(1\otimes d)|_S$ を積で定める。
$T=\{(d,e):\lambda d+\rho e=0\}$ と置き、
\[
 E_0=D,\qquad E_{n+1}=\operatorname{pr}_1(T\cap E_n^2)
                 \cap\operatorname{pr}_2(T\cap E_n^2)
\]
を安定するまで計算して $E$ とする。$L_v(h)=vh$, $R_v(h)=hv$ と書けば
\[
                 P_B=E/(L_H+R_H)
\]
は線型空間である。$B$ の同一視と点列を指定した上のOre表示の同型類は
$P_B$ と全単射である。$B$ の同一視を外すと $\Aut_{\mathrm{gr}}(B)$ の作用で割る。
点列も外した有限四辺の枝では、標的の高々二つの指定点でこの条件を試す。
$T\cap E^2$ は $P_B$ 上に一つの線型自己同型 $F_B$ のグラフを誘導し、
添字を一つ移すとパラメータは $F_B$ で移る。
\end{theorem}
\begin{proof}
各段階で計算するのは部分空間の交叉と射影であり、真に減る回数は高々18である。
両方向に延びる列の各項は常に $E_n$ に入る。逆に安定後は
$T\cap E^2$ の両射影が $E$ に全射で、任意の点を前後に延長できる。
式\eqref{eq:ternary-shift-derivation}により、これが全ての実現可能な初期値である。
$L_{v_i}-R_{v_{i+1}}$ は任意の列 $v_i\in H$ に対して同式を満たし、
従って $L_H+R_H\subset E$ である。

さらに
\[
                   \ker\rho=R_H,\qquad \ker\lambda=L_H .
\]
実際、$\rho d=0$ は生成元上で $h\mapsto d(h)$ とした次数1の
左 $B$ 加群写像 $B_{\ge1}\to B$ が定まることと同値である。
左イデアルの表示の最初の関係は $S$ が生成するためである。
$0\to B_{\ge1}\to B\to k\to0$ と左右AS条件
$\operatorname{Hom}_B(k,B)=\operatorname{Ext}^1_B(k,B)=0$ により、
この写像は $B\to B$ に一意に延びるので右乗法である。反対側も同様である。

二つの列の初期項が $P_B$ で一致すれば、位置0,1の $v_0,v_1$ で補正して
初期項を一致させる。次の差は $\ker\rho$ に入り $v_2$ だけで消せる。
順方向にこれを繰り返し、逆方向は $\ker\lambda$ と $v_{-1},v_{-2},\ldots$ を使う。
従って一位置の一致で全位置の同型が実現する。必要性は前系による。
この議論は指定した点列と $B$ を保つ場合には辺の有限性を使わない。
同じ初期類からの全ての次項が同じ類を与え、直前の項にも同じことが成立するので、
線型関係の像は両方向に一価、すなわち線型自己同型のグラフになる。

有限四辺の条件を課すには、$E$ 内で次項を三回取って四つの辺の六双一次式の
零点スキームの有限性を調べればよい。実現した対象は標準ASで四周期であり、
初期項が等しい全ての延長は指定点を保って同型なので、この検査は選択によらない。
指定点を保つ最小周期は $j=1,2,4$ の順に
$F_B^j[d]$ が $[d]$ と同じ自己同型軌道にあるかで得られる。
有限四辺では指定点を標準四移動が固定することをすでに示した。
指定点を外す場合は同じ移動先の他の三次元接点も候補に入れる。
\end{proof}

上記有理例の三生成部分では $E_n$ の次元は $18,15,15$、$L_H+R_H$ は六次元である。
九次元の商で $F_B^2=1$ となり、固有値 $+1,-1$ の重複度は順に $3,6$ である。
反可換三生成基底でも同じ次元と重複度を得る。可換三生成基底では
$E=D$、補生成元の取り替えによる部分空間は三次元で、十五次元の商上の移動は恒等である。
ここではすべて $X_i h-hX_{i+1}$ の符号規約を使い、
$B$ の自己同型群による商を取る前の数値を記している。
